\section*{Chapter \ref{ch:iv2:vol-surface-and-vol-trading} --- The Vol Surface and Vol Trading}
\markboth{\protect\omTocMark{}{Answers: The Vol Surface and Vol Trading}}{\protect\omTocMark{}{Answers: The Vol Surface and Vol Trading}}

\begin{solution}{iq:iv2:vol-surface-and-vol-trading:1}
It is about 22.4\%.

The squares of the returns add to 10, in units of squared percent. So the daily variance is $10/5 = 2$, and the daily volatility is about 1.414\%. Multiply by $\sqrt{252}$.

Daily means are so small compared with daily volatility that estimating the mean adds noise and changes little. A mean of 5 basis points is one twenty-fifth of a daily volatility of 1.25\%.

\iqwrong{Multiplying by 252. Variances scale with time, and volatilities with its root.}
\end{solution}

\begin{solution}{iq:iv2:vol-surface-and-vol-trading:2}
About 1.5\% a day. The 16 is $\sqrt{252}$, near 15.9. With every calendar day, it is $\sqrt{365} \approx 19$, which gives about 1.26\%.

\iqwrong{Treating 1.5\% as the expected move. It is one standard deviation. The expected absolute move is 0.8 times that, about 1.2\%.}
\end{solution}

\begin{solution}{iq:iv2:vol-surface-and-vol-trading:3}
You need about \$1.57, which is one daily standard deviation. In the model, the stock beats that about 32\% of the time.

A day's gain is $\tfrac12\Gamma m^2$ and a day's theta is $\tfrac12\Gamma S^2\sigma^2\Delta t$. They are equal when $m = S\sigma\sqrt{\Delta t} = 100 \times 0.25/\sqrt{252} \approx 1.57$. So you lose on about two days out of three, in small amounts, and win on one in three, in larger ones.

\iqwrong{Expecting to make money more than half the time. A long gamma position loses on most days, even when it is priced fairly.}
\end{solution}

\begin{solution}{iq:iv2:vol-surface-and-vol-trading:4}
It is about 18.5\%. Variance, weighted by time, adds up across independent periods, and volatility does not.

$0.22^2 \times 0.5 = 0.0242$, and $0.25^2 \times 0.25 = 0.0156$. The forward variance is $(0.0242 - 0.0156)/0.25 = 0.0343$. Its root is 18.5\%.

\iqwrong{Giving 19\%, from $2 \times 22 - 25$. That averages volatilities, and it is close here but gets further off for a steeper curve.}
\end{solution}

\begin{solution}{iq:iv2:vol-surface-and-vol-trading:5}
Total variance falls with time, and that is impossible. It would imply a negative forward variance. Sell the three-month option and buy the six-month one, which pays you about \$0.34.

Total variance is $0.09 \times 0.25 = 0.0225$ at three months, and $0.04 \times 0.5 = 0.02$ at six. The three-month call is about 5.98, and the six-month about 5.64. With zero rates and no dividends, the longer one can only be worth more. So at the first expiry, you can always sell the longer call for more than you owe on the shorter one.

\iqwrong{Saying short-dated volatility can be higher than long-dated. It can be, and it often is. Here it is too high even in total variance.}
\end{solution}

\begin{solution}{iq:iv2:vol-surface-and-vol-trading:6}
Lead with this: index returns are not lognormal. Big falls are likelier than big rises, and volatility rises in a fall. Second, demand: investors want puts to insure their holdings, and dealers who sell them charge for the risk. Third, correlation rises in a fall, so index volatility goes up even more than stock volatility.

The flat model has one constant volatility. A skew says the risk-neutral distribution has a fat left tail. Historically, index skews became steep only after the 1987 crash.

To build your own answer, start from the shape of the distribution, then give the supply and demand behind it, and finally say what you would test. For instance, you could ask whether the skew is steeper for an index than for its members, which it is.

\iqwrong{Saying that the market disagrees with Black-Scholes because the model is wrong, and stopping there. The interviewer wants to know in what way it is wrong, and who is paying for what.}
\end{solution}

\begin{solution}{iq:iv2:vol-surface-and-vol-trading:7}
The variance notional is \$2,500 per variance point. You make \$562,500 at 25, and you lose \$437,500 at 15. The gain is larger because you are paid on the square of volatility.

The notional is $100{,}000/(2 \times 20)$. The payoff is $2{,}500(\sigma^2 - 400)$. So $625 - 400 = 225$, and $225 - 400 = -175$. Around the strike, a point of volatility is worth \$100,000, which is why it is called vega notional.

\iqwrong{Treating the payoff as linear in volatility. That is a volatility swap. A variance swap is convex, and a seller feels it in a crash.}
\end{solution}

\begin{solution}{iq:iv2:vol-surface-and-vol-trading:8}
It is about 3.2 points, and you need 200 days for 1 point.

The rule is $\sigma/\sqrt{2n}$. So $20/\sqrt{40} = 3.16$. For 1 point, $2n = 400$. Over 20 days, a true 20\% looks like anything between about 14 and 26.

\iqwrong{Comparing a 20-day realised number to implied and calling a 2-point gap a signal. A gap of that size is inside the noise.}
\end{solution}

\begin{solution}{iq:iv2:vol-surface-and-vol-trading:9}
It is 36\%. For a large index with equal volatilities and pairwise correlation $\rho$, $\sigma_I^2 \approx \rho\sigma^2$. So $\rho = (18/30)^2$.

A rise means the options market is charging more for stocks moving together, or that index options have got dearer against stock options. It is often a sign of fear, though you should check that demand for index protection has not simply risen.

\iqwrong{Using $18/30 = 60\%$. Correlation multiplies variance, and it needs the square.}
\end{solution}

\begin{solution}{iq:iv2:vol-surface-and-vol-trading:10}
At 100, the implied volatility is about 19.98\%. At 80, it is about 22.91\%, and at 130, about 23.60\%. So uncertainty about volatility creates a smile. At the money, you get about the average because an at-the-money price is almost linear in volatility.

The average of the two volatilities is 20\%, and the root mean square is 22.36\%. For a far strike, price is convex in volatility, so the mixture is priced above the average, and its implied volatility is higher. That is what gives a smile with a vol of vol and no skew: a symmetric mix gives a symmetric shape in log strike, as long as there is no correlation with the stock.

\iqwrong{Thinking a smile proves the model has jumps. A random volatility alone can do it, so a smile does not tell you which story is true.}
\end{solution}

\begin{solution}{iq:iv2:vol-surface-and-vol-trading:11}
It is about 0.0196 per dollar from these rounded prices, close to the true 0.01985. A negative answer would be an arbitrage.

The butterfly gives $(C_{95} - 2C_{100} + C_{105})/5^2$. The density is the second derivative of the call in strike\index{Breeden-Litzenberger}. The lognormal density at $K$ is $\varphi(d_2)/(K\sigma\sqrt T)$, here $\varphi(-0.1)/20$.

Rounding matters here. With prices rounded to the cent, the sum has an error of up to 2 cents, or 0.0008 in the density, which is 4\% of the answer. With unrounded prices, the butterfly gives 0.0198.

\iqwrong{Using strikes too close together on rounded quotes. Narrower butterflies pick up more rounding, and the noise can make the density look negative when it is not.}
\end{solution}

\begin{solution}{iq:iv2:vol-surface-and-vol-trading:12}
Sticky strike: the at-the-money is 19.5\%, and the 100 strike stays at 20\%. Sticky moneyness: the at-the-money stays at 20\%, and the 100 strike goes to 20.5\%. Your profit is about 20 cents higher under sticky moneyness.

Under sticky strike, each strike keeps its own volatility, so the new at-the-money is the old 105 strike: $20 - 0.5$. Under sticky moneyness, the whole smile moves with the stock, so the 100 strike is now 5\% below the money: $20 + 0.5$. The gap is half a point, at 0.39 a point, which is about 0.2.

\iqwrong{Saying one of these is true. They are two ends of what the market does, and which one is closer depends on the size and speed of the move.}
\end{solution}

\begin{solution}{iq:iv2:vol-surface-and-vol-trading:13}
It is 0.54 under sticky strike, about 0.62 under sticky moneyness and about 0.46 under the local volatility model. None is right in all markets. So measure the empirical move of the surface with the stock, and hedge accordingly.

The general rule is $\Delta = \Delta_{\text{BS}} + \text{vega}\,\partial\sigma/\partial S$. Under sticky strike, the last term is zero. Under sticky moneyness, $\sigma$ depends on $K/S$, so a rise in $S$ lifts the volatility at a fixed strike by 0.002 per dollar. The extra term is $39.7 \times 0.002 = 0.079$. Local volatility moves the other way: $0.54 - 0.079$.

\iqwrong{Quoting one delta and thinking it has no model in it. A delta holds a view on how the smile moves.}
\end{solution}

\begin{solution}{iq:iv2:vol-surface-and-vol-trading:14}
One point of the 1.5 is the slide along the smile, and half a point is the smile itself moving up. The option at the old at-the-money strike is 2\% above the new at-the-money level, so its volatility went from the old at-the-money to about half a point above it.

The new at-the-money is the strike that was 2\% below, so sliding gives $2 \times 0.5 = 1$. The remaining 0.5 is a parallel shift. The old strike moves by just that parallel shift, as it stays where it was along the original curve.

\iqwrong{Saying the implied volatility of your option went up 1.5 points. It did at the new at-the-money strike, and your option's strike is not the new at-the-money one.}
\end{solution}

\begin{solution}{iq:iv2:vol-surface-and-vol-trading:15}
You make about \$363,600. That is about 3.6 points of vega. Implied may stay above realised because investors pay to be protected from volatility rising in a fall, and sellers want to be paid for that risk. The catch is that the seller loses in large amounts at the worst time.

The payoff is $100{,}000/(2 \times 22) \times (484 - 324) = 2{,}273 \times 160$. The strike was four points above realised, but convexity makes it worth a little less than four.

On the evidence, delta-hedged index options have lost money on average, a result known as a negative volatility risk premium. On a February day in 2018, the S\&P 500 fell 4\% and its volatility index rose 20 points. So that gap does not come free.

\iqwrong{Treating the average gap as a free lunch. Whether it pays is a matter of size and of the tail.}
\end{solution}

\begin{solution}{iq:iv2:vol-surface-and-vol-trading:16}
The estimator is Parkinson's\index{Parkinson estimator}: the squared log range, times $1/(4\ln 2) \approx 0.36$. For a driftless Brownian motion, $E[\text{range}^2] = 4\ln 2\,\sigma^2$. That gives about one-fifth of the variance of a squared return, so about five days' worth of close-to-close data from one day.

The biases: real prices are seen only at discrete times, and the true high is never seen, which lowers the estimate. The overnight gap is left out entirely. Bid-ask bounce raises the range.

So use a form that adds an overnight term, and compare with a realised volatility from intraday returns.

\iqwrong{Using it without the factor. The squared range has a mean about 2.77 times the variance.}
\end{solution}

\begin{solution}{iq:iv2:vol-surface-and-vol-trading:17}
It shows about 17.4\%. It looks low since the screen counts three days of time, while the stock is only open for one. On Monday morning, the same screen shows about 30\%.

Equating the variance, $\sigma_c^2 \times 3/365 = 0.25^2/252$. So $\sigma_c = 25\% \times \sqrt{365/756} \approx 17.4\%$.

On Monday, one trading day is one calendar day. Then $25\% \times \sqrt{365/252} \approx 30\%$. So the very short end of a calendar-time curve has a weekly sawtooth, and it means nothing by itself. A trader who does not see this will find false trades.

\iqwrong{Buying Friday's volatility because it looks cheap, or selling Monday's because it looks dear. Use a clock that gives trading days and event days their real weight.}
\end{solution}

\begin{solution}{iq:iv2:vol-surface-and-vol-trading:18}
The base volatility is about 33.2\%, and the event's implied move is about 7.0\%, so the expected absolute move is about 5.6\%. You assume that the base volatility is constant, and that the event is one independent return.

The total variance is $0.60^2 \times 5/252 = 0.00714$ for the short option, and $0.40^2 \times 25/252 = 0.01587$ for the long one. Their difference is base variance over 20 days. So the base is $0.00873 \times 252/20 = 0.110$, whose root is 33.2\%. The event variance is $0.00714 - 0.110 \times 5/252 = 0.00496$, whose root is 7.04\%. The absolute move has the factor $\sqrt{2/\pi} = 0.8$.

\iqwrong{Using the 60\% as the event's volatility. Most of the short option's variance comes from the event, but the base days are in it too.}
\end{solution}

\begin{solution}{iq:iv2:vol-surface-and-vol-trading:19}
The implied volatility is about 32.6\%. After the event, the straddle is worth about 5.62 if unchanged, 6.20 at 103, 7.66 at 106 and 10.60 at 110. You break even at a move of about \$5.7.

Total variance adds: $0.25^2 + 0.06^2/(30/365)$, whose root is 32.6\%. The straddle costs 7.46 at that level. After one day, its volatility falls to 25\%, and that takes away about \$1.83 even with no move: the vol crush.

So a straddle is a bet on a move larger than the market's expected absolute move, and that is 4.8\%. You need nearly 6\%. Because of the remaining value of the option after the event, the break-even is smaller than the straddle's price of 7.46.

\iqwrong{Comparing the expected move with the straddle price alone. The straddle still has value after the event, so you have to price what is left.}
\end{solution}

\begin{solution}{iq:iv2:vol-surface-and-vol-trading:20}
The fair variance strike is $\sqrt{0.04 + 0.0025} \approx 20.6$. The fair volatility strike is 20. A volatility swap pays less in a big move, as it pays on volatility, which is linear in its own value, and not on its square.

The variance strike is the root of the mean of the variance: $E[\sigma^2] = 20^2 + 5^2 = 425$. The volatility swap is just the mean of the volatility.

It is hard to hedge because the square root is concave in variance. A static strip of options replicates variance, not volatility. So a volatility swap needs a model of vol of vol, and dynamic trading in variance. The gap between the two strikes, here about 0.6, is the convexity adjustment\index{convexity adjustment}.

\iqwrong{Setting the two strikes equal. Doing this ignores the variance of volatility, and that is the exact thing a variance swap charges for.}
\end{solution}

\begin{solution}{iq:iv2:vol-surface-and-vol-trading:21}
The index's volatility is about 16.8\%. If correlation realises at 0.2, the index's volatility is about 13.9\%, and you gain 88 variance points on the index with no change on the stocks. At a stock volatility of 33\%, the index is at about 15.3\%, and you also gain on the stocks. This is called dispersion\index{dispersion trading}. You are short correlation, and that hurts in a crash, when correlation jumps.

The formula is $\sigma_I^2 = \sigma^2[1/n + (1 - 1/n)\rho]$. With $n = 50$, $\sigma = 30\%$ and $\rho = 0.3$, it is 0.02826. At $\rho = 0.2$, it is 0.01944, so 88 points less. The same formula at 33\% gives 15.3\%.

You collect the implied correlation less the realised, which is a risk premium, and you pay for it in a crisis. It is also expensive to trade, as you need 50 stock positions, and they all need hedges.

\iqwrong{Calling it a pure volatility trade. It is mainly short correlation. With weights chosen to be vega-flat, you are still exposed to volatility levels.}
\end{solution}

\begin{solution}{iq:iv2:vol-surface-and-vol-trading:22}
With skew, it costs about \$2.69. At flat 20\%, it would cost about \$4.38. It is cheaper since the put you sell has a higher volatility and so a higher price: about 5.27, against 3.59.

The 100 put at 20\% costs 7.97. So $7.97 - 5.27 = 2.69$ and $7.97 - 3.59 = 4.38$.

Whoever buys a downside put pays for skew, and someone who sells a further put recovers much of it. That is why an investor who wants cheap protection buys a put spread, and gives up the tail.

\iqwrong{Using one volatility for all strikes in a spread. The spread is the part of the smile where the skew matters most.}
\end{solution}

\begin{solution}{iq:iv2:vol-surface-and-vol-trading:23}
It does not meet it. The condition is $2\kappa\theta \geq \xi^2$, and here $0.16 < 0.25$. So the variance can reach zero, and a simple simulation will send it negative.

If you simulate it by Euler, use full truncation: you put $\max(v, 0)$ into the drift and the diffusion. Or use the quadratic-exponential scheme, which matches the exact moments better. Most calibrated models fail the condition, so you cannot assume it away.

\iqwrong{Forcing the condition in the calibration. That changes the smile you can fit, and it is not needed for the price to exist.}
\end{solution}

\begin{solution}{iq:iv2:vol-surface-and-vol-trading:24}
Interpolate linearly in total variance, in the clock you trust. You get about 21.5\%.

The totals are $0.04/12 = 0.00333$ and $0.0484 \times 0.25 = 0.0121$. Halfway in time gives 0.00772, for two months. Multiply by 6 and take the root.

Linear interpolation in volatility gives 21\%, and it can yield falling total variance on a steep curve. Linear in total variance cannot, as long as the input is itself increasing. So check that total variance at each strike, at a fixed moneyness, rises with expiry. With events, you weight the clock, and you do not go straight across them.

\iqwrong{Treating it as a question about one number. Across a whole surface, the same rule must hold at each moneyness, or calendar arbitrage comes back in the wings.}
\end{solution}

\begin{solution}{iq:iv2:vol-surface-and-vol-trading:25}
It is worth about \$386,250. The mark is simple since variance adds over time. Its vega has fallen to about half of what it was, because only half the time is left.

The expected total variance is $\tfrac12 \times 625 + \tfrac12 \times 484 = 554.5$. So the value is $2{,}500 \times (554.5 - 400) = 386{,}250$, at zero rates.

The vega on a remaining-variance point is $2{,}500 \times \tfrac12$, per point of variance. At a volatility of 22, a point of volatility is $2 \times 22$ variance points. So the vega is about \$55,000, against \$100,000 at the start. It is linear in the remaining time, so a variance swap has no theta surprises, and it is a clean exposure.

\iqwrong{Marking on the current implied alone. Most of a seasoned variance swap's value comes from what is already realised.}
\end{solution}

\begin{solution}{iq:iv2:vol-surface-and-vol-trading:26}
A strip weighted by $1/K^2$ has the same dollar gamma at any stock price. So its gain per unit of variance does not depend on where the stock goes. With equal weights, dollar gamma rises with the stock price. The fair strike here is 20\%, as a flat smile has nothing else to give.

A hedged option earns $\tfrac12S^2\Gamma(\sigma_r^2 - \sigma_i^2)\,dt$. So for a clean variance exposure, $S^2\Gamma$ must be constant. Each option's gamma depends on $S/K$ in a way that makes this work exactly when $f^{\prime\prime}(K) \propto 1/K^2$.

In a numerical check at stock prices of 80, 100 and 120, using strikes from 10 to 1,000, the strip's dollar gamma stays the same to within 0.2\%. The one with equal weights rises with the square of the stock price.

The strip is the second-derivative expansion of $-2\ln S$. That is why this is the log contract\index{log contract} in disguise, and that is why the fair strike is $-(2/T)E[\ln(S_T/S_0)]$.

\iqwrong{Hedging a variance swap with an at-the-money straddle. It has the right exposure only at one stock price, and the exposure fades as the stock moves.}
\end{solution}

\begin{solution}{iq:iv2:vol-surface-and-vol-trading:27}
Such a smile cannot be right, because a call would rise in value at higher strikes. The largest slope allowed is 2 in total variance per unit of log strike. For a stock at 100, at a strike of 400, the highest implied volatility is about 167\%.

Roger Lee's\index{Lee's moment formula} moment formula says that total variance can grow at most like $2|k|$ for large $|k|$, where $k = \ln(K/S)$. Above that, the Black-Scholes price with that total variance does not fall to zero as $K$ grows. A call's price must fall with strike. And for a large $K$, it must go to zero.

In a check, a total variance of $3|k|$ makes the call's value rise with strike all the way out, from 42 at a strike of 300 to 82 at $10^8$. With $1.5|k|$, it falls far out, from 23 at 1,000 to 17 at $10^8$, as it should.

For the number: $k = \ln 4 = 1.386$, so the total variance is at most 2.77, and its root is 1.665 at one year.

\iqwrong{Extrapolating a polynomial in the wings. Fitting it through the quotes makes the prices of far options nonsense, and the error then falls into the variance swap strike, which integrates those options.}
\end{solution}

\begin{solution}{iq:iv2:vol-surface-and-vol-trading:28}
The formula is $\sigma_{\text{loc}}^2(K, T) = 2\partial_TC/(K^2\partial_{KK}C)$. For a flat surface, you get back 20\%. For the mixture at a strike of 100, you get about 17.3\%. It is so different because it is a conditional average, while the implied volatility is an average over paths.

For the flat case, $\partial_TC = S\varphi(d_1)\sigma/(2\sqrt T)$, and $K^2\partial_{KK}C = K\varphi(d_2)/(\sigma\sqrt T)$. When you divide them, you get $\sigma^2$.

Dupire's\index{Dupire's formula} formula is equal to the expected instantaneous variance, given that the stock is at that level at that time. In the mixture, a stock at 100 after a year is much likelier to be in the 10\% state, since its density there is taller. So the weights are about 75\% and 25\%, and $\sqrt{0.75 \times 0.01 + 0.25 \times 0.09} \approx 17.3\%$. Implied volatility at the money is about the mean of the two states.

\iqwrong{Reading local volatility as a forecast of future realised volatility. It is what fits today's option prices. It is not an estimate of what the stock will do.}
\end{solution}

\begin{solution}{iq:iv2:vol-surface-and-vol-trading:29}
The implied volatilities are about 22.05\%, 19.98\% and 18.13\%, by simulation. That is a slope of about $-0.1$ per unit of log strike, half of the local slope. The rule is that the implied slope is half of the local slope for a short maturity.

Implied volatility at a strike is close to an average of the local volatility over the likely path from spot to that strike. If that path runs in a straight line in log price, the average is halfway. So the implied slope is half of the local one.

The test is a simulation under the model. Its slope is within 0.01 of the rule, even at one year.

\iqwrong{Fitting a local volatility of the same slope as the observed skew. The fit is then wrong by a factor of two, and the model's vanilla prices are off.}
\end{solution}

\begin{solution}{iq:iv2:vol-surface-and-vol-trading:30}
The at-the-money fell by about 0.97 points, which is about twice the skew's own 0.49 over that move. And the 100 strike's volatility fell by half a point, so the smile moved to the left as the stock moved to the right. Markets tend to do the opposite, moving the smile with the stock.

So the model's delta has the wrong correction. It subtracts a vega term where the market needs one added. Its hedges are therefore worse than Black-Scholes in many tests, which is the well-known critique by Hagan and co-authors.

Local volatility is still useful. It reprices the vanilla options exactly, and it gives barrier-like products a value consistent with every vanilla price. But the dynamics are something you must add yourself.

\iqwrong{Thinking that a model that fits all options must hedge them well. A fit is a static statement, and a hedge is about how the fit changes.}
\end{solution}

\begin{solution}{iq:iv2:vol-surface-and-vol-trading:31}
The structure is about 22.5\%, 26.2\% and 28.9\%, rising toward $\sqrt{0.09} = 30\%$. The correlation sets the slope, and the vol of vol sets the curvature. The at-the-money is below 20\%, because at-the-money prices are nearly linear in volatility and the mean of a volatility is below the root of its mean variance.

The expected average variance over $T$ is $\theta + (v_0 - \theta)(1 - e^{-\kappa T})/(\kappa T)$. So at one year, $0.09 - 0.05 \times 0.432 = 0.0684$, and the root is 26.2\%. A simulation confirms the 26.2\%.

For the smile: with a correlation of $-0.7$, a falling stock raises volatility, which gives a fatter left tail, and that is the skew. With zero correlation, the smile is symmetric in log strike. Its height is the vol of vol.

\iqwrong{Thinking the mean-reversion speed sets the skew. It governs how fast the skew fades with maturity.}
\end{solution}

\begin{solution}{iq:iv2:vol-surface-and-vol-trading:32}
At the money, it is about 20.11\%. At 80, it is about 21.93\%, and at 120, about 19.44\%. A simulation gives 20.14, 21.87 and 19.44. Its advantage is that the smile moves in the right direction with the stock. Its weakness is that its expansion can give negative densities at low strikes for long maturities.

The at-the-money formula, with $\beta = 1$, is $\alpha[1 + (\rho\nu\alpha/4 + (2 - 3\rho^2)\nu^2/24)T]$. So $0.2 \times (1 - 0.006 + 0.01153) = 0.2011$, or 20.11\%. Away from the money, the formula uses $z = (\nu/\alpha)\ln(F/K)$.

It is a formula, and not a model solution, so it can be wrong in low strikes and long maturities. That matters most in rates, where it is the market's quoting tool. Check it by simulation, or against a density computed from the formula's own prices.

\iqwrong{Using SABR as a pricing model for path-dependent products. It is a good interpolator and it gives a delta. It does not fit the term structure as a whole.}
\end{solution}

\begin{solution}{iq:iv2:vol-surface-and-vol-trading:33}
At one month, the implied volatilities at 90, 100 and 110 are about 20.8\%, 15.6\% and 15.2\%. At one year, they are about 16.7\%, 16.3\% and 16.0\%. The skew is so steep short-dated, because a diffusion hardly moves at all in a month, so only the jump can reach 90. A pure diffusion's skew flattens out at short maturities unless its volatility is extreme near those strikes.

The call is $\sum_jp_jC_{\text{BS}}(F_j, K, \sigma, T)$, where $p_j$ are Poisson with mean $\lambda T$, and $F_j = 100e^{-\lambda JT}(1 + J)^j$ with $J = -0.2$. The factor $e^{-\lambda JT}$ puts back the expected loss from the jumps, so that the stock is a martingale.

So in the short end, the skew is mostly about the jump, with a price set by its chance and size. Over a year, the jump gets averaged into the mix, and the skew flattens quickly. Local volatility can fit this surface today, but only with extreme volatility levels near expiry.

\iqwrong{Trying to fit a steep one-month skew with diffusive stochastic volatility alone. It needs an extreme vol of vol or very quick mean reversion, and then it breaks at longer maturities.}
\end{solution}

\begin{solution}{iq:iv2:vol-surface-and-vol-trading:34}
A down jump of 20\% has a squared log return of 0.0498, and the strip gives 0.0463, so you lose 0.0035 per unit of variance notional. An up jump of 20\% has a squared log return of 0.0332, and the strip gives 0.0354, so you gain 0.0021.

The strip pays $2(r - \ln(1 + r))$ on a simple return $r$, as it is the log contract less a delta hedge. That matches $r^2$ only for small moves. For $r = -0.2$, $2(-0.2 + 0.2231) = 0.0463$. For $r = +0.2$, it is $2(0.2 - 0.1823) = 0.0354$.

A short swap with a strip is therefore short a down jump. So a variance swap's fair strike is above what continuous replication would suggest, with a negative-skew distribution. It is also why swaps on single stocks are often capped: the seller's loss in a default is huge.

\iqwrong{Treating a variance swap's replication as exact. It is exact only for a continuous path. With a jump, you have the term $r^3$, and its sign does the damage.}
\end{solution}

\begin{solution}{iq:iv2:vol-surface-and-vol-trading:35}
The fair strike from the strip is about 20.54\%, in variance 0.0422. The rule gives $0.04(1 + 0.03) = 0.0412$, or 20.30\%. The strike is above the at-the-money because the strip includes the downside puts, and with a skew those are dear.

The strip is $(2/T)\int(\text{OTM option})/K^2\,dK$, and it covers every strike. Its weights are heavier on the downside, and the put wing is where the smile is highest. The variance is therefore above the at-the-money, whatever the direction of the skew, and the rule's $b^2$ shows that.

The rule is rough: it is a second-order expansion with a linear skew, and it leaves out the curvature. So it is low by 0.24 points here. It is a useful check, and not a price.

\iqwrong{Quoting at-the-money as the variance strike. In an index market, it can be several points below.}
\end{solution}

\begin{solution}{iq:iv2:vol-surface-and-vol-trading:36}
The index has a volatility of about 13.4\% in the normal state and about 26.8\% in the other. With the same forward, the smile is about 17.4, 16.1 and 17.0, at 80, 100 and 120. That is a smile, and almost no skew. To get a skew, you must make the high-correlation state one in which prices have fallen.

For a large index, $\sigma_I^2 = \rho\sigma^2$. The mixture's weights are 0.8 and 0.2. A check with the high-correlation state's forward at 85 and the other at 103.75, which keeps the mean at 100, gives a marked downward skew.

Index skews are steeper because a fall in the index tends to come with higher correlation, on top of higher stock volatility. A single stock has its own news, which can be good or bad. As a documented result, single-stock risk-neutral distributions are much less negatively skewed than the market's.

\iqwrong{Saying that high correlation by itself causes the skew. It adds to volatility in one state. The link between that state and falling prices is what makes it a skew.}
\end{solution}

\begin{solution}{iq:iv2:vol-surface-and-vol-trading:37}
The mean is about 0.93, on a premium of 4.61, and the standard deviation is about 0.75. You lose about 9\% of the time, and the 1st percentile is about $-0.68$. So the Sharpe ratio per trade is about 1.2, and this is an unrealistically good case.

The mean is the difference of the two prices, at 20\% and at 16\%, and that is 0.92. The vega route agrees: $4 \times 0.23$.

Your model has no jumps, and its volatility is known. In reality, realised volatility is uncertain, and a big move can cost many months of gains. So you size from a stress: for example, a 20\% gap on the straddle costs nearly \$15.4 here, or about 17 times the mean profit. Then you decide how much of that you can bear.

\iqwrong{Sizing with the model's standard deviation. A short-gamma trade looks best in the statistics in exactly the period before it fails.}
\end{solution}

\begin{solution}{iq:iv2:vol-surface-and-vol-trading:38}
By Ito, $\tfrac12\int\sigma^2\,dt = \int dS/S - \ln(S_T/S_0)$. So the fair variance is $-(2/T)E[\ln(S_T/S_0)]$, and the Carr-Madan expansion of the log gives
\[
K_{\text{var}} = \frac2T\Big[\int_0^{S_0}\frac{P(K)}{K^2}\,dK + \int_{S_0}^\infty\frac{C(K)}{K^2}\,dK\Big].
\]
The exact value is 0.04218, or 20.54\%. With strikes from 50 to 200, a step of \$1 gives 0.04205, \$5 gives 0.04250 and \$10 gives 0.04378. From 70 to 140 in steps of \$5, it gives 0.04096.

The two errors differ. A coarse step overstates, since the integrand is convex, as with the trapezoid rule on a kink at $S_0$. That is the 0.04250 and 0.04378. A truncated range understates, as it leaves out options, and that is 0.04096 and part of the 0.04205.

In a stress, truncation matters more. When a stock falls, more of the variance is in strikes outside the original range. That is why a replicating book's gamma falls off beyond its last strike. The difference between a listed and an actual strip is what a fund sees when it thinks it has hedged and has not.

\iqwrong{Trying only finer strikes. The wing's cost grows with volatility, while the step's error is easy to reduce. A wider range is where most of the money is.}
\end{solution}

\begin{solution}{iq:iv2:vol-surface-and-vol-trading:39}
Condition on the path of volatility. Then the stock is lognormal with a variance of $\bar v = \frac1T\int v\,dt$, so its call price is $C_{\text{BS}}(\sqrt{\bar v})$. The average over paths of that is the price. This is the Hull-White\index{Hull-White mixing} mixing formula. In a test, the two methods give the same smile to within noise, and the mixing estimate's standard error is about 5 times smaller at the money.

For the mixing result, simulate only the variance and record $\bar v$. A plain simulation also needs the stock's noise.

So the second is a conditional Monte Carlo estimate: it removes the stock's noise exactly. That is why it reduces variance, for free.

With correlation, conditioning on $v$ also changes the stock's mean, since the stock's noise is linked to that of the variance. Then the conditional stock price is lognormal with a different forward: you scale $S$ by $\exp(\rho\int\sqrt v\,dW - \tfrac12\rho^2\int v\,dt)$ and use the remaining $(1 - \rho^2)\bar v$. The Romano-Touzi form of the same idea handles that.

\iqwrong{Using a mixture with the wrong condition. If the correlation is not zero, the plain formula gives no skew at all, and that is not correct.}
\end{solution}

\begin{solution}{iq:iv2:vol-surface-and-vol-trading:40}
The local variance is $\sigma^2_{\text{loc}}(x, t) = w_1\sigma_1^2 + w_2\sigma_2^2$, where each $w_i$ is proportional to the density of state $i$ at the log price $x$ at time $t$. A plain daily simulation misses: at a strike of 70, it gives an implied volatility of about 23.2\% against 24.97\%. A partial differential equation on a time grid that bunches steps near zero solves it to within 0.02 of a point.

The formula comes from Gyöngy's\index{Gyöngy's theorem} theorem: the local variance is $E[v \mid S_t = S]$. The state's weight is its share of the density. Using this, one model with a single random factor matches all of the mixture's marginal distributions.

Near time zero, the local variance is singular. The two densities have very different widths, so at small $t$, the 30\% state takes over very close in, and $\sigma_{\text{loc}}$ swings from 10\% to 30\% over a very small distance. A daily step jumps over this. So the path misses the early spread, and the far strikes come out too low. The expected price also drifts from 100.

The solution is to take steps that are small at the start, such as $t_i = (i/m)^2$, and to use an implicit grid in $x$. Since local volatility repricing should hold exactly, this is a useful check of any calibration.

\iqwrong{Blaming the sampling error for this. With more paths, the result stays wrong by the same amount. Only more steps early on will fix it.}
\end{solution}
